\documentclass[11pt,a4paper]{article}
\usepackage[margin=22mm]{geometry}
\usepackage{amsmath,amssymb,amsthm,booktabs,array}
\usepackage[colorlinks=true,urlcolor=blue,linkcolor=black,citecolor=black]{hyperref}
\usepackage{microtype}
\setlength{\parskip}{0.5em}
\setlength{\parindent}{0pt}
\newcommand{\T}{\mathcal T}
\newtheorem*{lemma}{Lemma}

\title{A bounded quantitative claim for the I01 follow-up (revised)\\
\large Charge-sector path fingerprints in a constrained phase interferometer}
\author{Working research record}
\date{13 November 2026\\[2pt]\small Supersedes the claim of 6 October 2026, whose record is kept unchanged.}

\begin{document}
\maketitle
\vspace{-1.2em}

\section*{Decision in one sentence}
Our working contribution is an \emph{exact, observable-level calculation for one specified digital circuit}, together with preregistered tests of its predictions. The circuit consists of a $P=0$ preparation and return, two transverse access pulses around a diagonal phase layer, an energy-correlated clock projection, and a phase-label allocation. The underlying mechanism is a coherent excursion out of a symmetry sector whose second-order return is sensitive to symmetry phases, and which symmetry averaging suppresses. That mechanism is established in the literature on gauge-violating errors and symmetry protection \cite{halimeh2020,lamm2020,tran2021}. We claim no priority for it, no derivation of three preferred spatial dimensions, and no usable fault-localization protocol.

\section*{What changed since 6 October 2026}
\begin{enumerate}\itemsep0pt
\item \textbf{Prior work.} The gauge-violation and symmetry-protection literature is now included, and the mechanism is no longer part of the claim.
\item \textbf{Phase labels.} A uniform label average is a $\mathbb Z_L$ twirl of the $U(1)$ generated by the charge $P$. The short-label failures of prevalidation stage~1 are aliasing; stage~27a reproduces them exactly.
\item \textbf{Path counts.} The counts are trinomial coefficients. The first opening at $n=3$ is a charge-balance statement (one spectator axis is needed), not a dimensional one.
\item \textbf{Composition of the contrast.} The measured contrast is dominated by interference, and a Pauli twirl of the pulses removes it (stage~27a).
\item \textbf{Noise.} The phase-damping null now follows from a lemma. The computed preparation- and readout-induced leakages are kept as numbers, not as a localization result.
\item \textbf{Scope.} Conjugate-leakage localization is removed from scope, because the shot budget is infeasible and repetition destroys the null.
\item \textbf{Encodings.} The encoding dependence found in stages 24--26 is added.
\end{enumerate}

\section*{Circuit and exact amplitude law}
For each of $n$ momentum axes, $k_i\in\{-1,0,1\}$ is encoded on a plus bit and a minus bit. Prepared and returned system states are supported on $S_n=\{\boldsymbol k:\sum_ik_i=0\}$. A source pulse $R_x^{(a)}(\varepsilon)$ precedes the diagonal phase layer $U(\boldsymbol\theta)$, and a return pulse $R_x^{(b)}(-\sigma\varepsilon)$ follows it, with $\sigma=\pm1$ ($\sigma=+1$ is the standard configuration). The clock is projected onto its zero outcome, and its phase is absorbed into the prepared amplitudes $|w_\tau\rangle$. With $A_0=\langle r|U|w_\tau\rangle$, $A_{ab}=\langle r|X_bUX_a|w_\tau\rangle$, $c=\cos\frac\varepsilon2$, $s=\sin\frac\varepsilon2$ and $c_n=\lceil\log_2(n+1)\rceil$ clock qubits,
\begin{equation}\label{eq:amp}
 p_\sigma=2^{-c_n}\bigl|c^2A_0+\sigma s^2A_{ab}\bigr|^2 .
\end{equation}
Odd-flip terms vanish because one bit flip cannot connect two $P=0$ support states of this encoding. Let $C_i$ and $D_i$ be the label-pooled block differences of $2^{-c_n}\mathrm{Re}(A_0^*A_{ab})$ and $2^{-c_n}|A_{ab}|^2$ between two common-phase blocks. Because $A_0$ is invariant under a common phase,
\begin{equation}\label{eq:angle}
 \Delta_{i,\sigma}(\varepsilon)=\sigma\,\tfrac12\sin^2\!\varepsilon\,C_i+\sin^4\!\tfrac\varepsilon2\,D_i,
 \qquad
 \Delta_i^{\rm tw}=\tfrac12(\Delta_{i,+}+\Delta_{i,-})=\sin^4\!\tfrac\varepsilon2\,D_i .
\end{equation}
Averaging over $\sigma$ is equivalent to a Pauli twirl of each pulse. For $\sigma=+1$, Eq.~\eqref{eq:angle} reproduces 36 Aer angle conditions within $1.16\times10^{-16}$ and the stage-15 $n=2,3,4$ contrasts within $1.03\times10^{-15}$. In stage~27a both signs, over 672 independent Aer circuits, obey Eq.~\eqref{eq:angle} within $1.3\times10^{-16}$, and the $\sigma$-average equals $\sin^4(\varepsilon/2)D_i$ within $3.6\times10^{-17}$.

\textbf{The contrast is interference.} At $n=3$ and $\varepsilon=\pi/2$ the twirled contrasts are $2.4\times10^{-8}$, $2.4\times10^{-9}$ and $4.3\times10^{-9}$ for the three allowed classes, against coherent contrasts of $2.04\times10^{-3}$, $5.47\times10^{-4}$ and $9.86\times10^{-4}$. The same-qubit value implies $D\approx9.7\times10^{-8}$, which is the small $\varepsilon=\pi$ residual found in stage~16 ($9.74\times10^{-8}$). It follows that randomized compiling \cite{wallman2016,hashim2021} applied to these pulses would remove the signal. Conversely, the twirl gives the finite-shot test a built-in negative control.

\section*{Path classes and the meaning of $n=3$}
Write $\T_m(q)=[z^q](z^{-1}+1+z)^m$ (OEIS A002426 and A005717 \cite{oeis}). The number of $P=0$ basis states whose two-flip path reaches the return support is:
\begin{center}\small
\begin{tabular}{@{}l l r r r r r@{}}
\toprule
Source--return class & Count & $n=2$ & $3$ & $4$ & $5$ & $6$\\
\midrule
Same qubit & $\T_n(0)$ & 3 & 7 & 19 & 51 & 141\\
Conjugate bit, same axis & $0$ & 0 & 0 & 0 & 0 & 0\\
Other axis, same sign & $2\T_{n-2}(-1)$ & 0 & 2 & 4 & 12 & 32\\
Other axis, opposite sign & $2\T_{n-2}(0)$ & 2 & 2 & 6 & 14 & 38\\
\bottomrule
\end{tabular}
\end{center}
The table was re-enumerated by brute force for $n=2$--$6$. The same-qubit count is simply $|S_n|$, because $X_aX_a=I$. The other-axis/same-sign path switches a $+1$ on one axis on and a $+1$ on another axis off. Its starting state is in $S_n$ only if the remaining $n-2$ axes carry total charge $-1$, which requires at least one spectator axis. The \emph{first opening at $n=3$ is therefore a charge-balance fact}; the class stays open for all $n\ge3$. A nonzero count is necessary, not sufficient, for a nonzero contrast.

\section*{Phase labels as a $\mathbb Z_L$ twirl}
Label $l$ multiplies every axis by $\theta_l=2\pi l/L$, that is, by $e^{i\theta_lP}\in\mathbb Z_L\subset U(1)$. A uniform label average keeps only charge differences $q=P(x)-P(y)\equiv0\pmod L$. For a common shift $\alpha$, the label-erased change is
\begin{equation}
 \delta_{\rm alias}=\sum_{q\equiv0\,(L),\,q\neq0}(e^{i\alpha q}-1)S_q,\qquad S_q=\sum_{P(x)-P(y)=q}a_x\bar a_y .
\end{equation}
For states with charges in $[-n,n]$, $L\ge2n+1$ removes every $q\neq0$. For $L=2n-1$ the components $q=\pm(2n-1)$ survive. This reproduces all 128 rows of the prevalidation-1 grid within $3.3\times10^{-16}$, including the 16 unrestricted short-label rows that responded to a common shift after label erasure.

\emph{Post hoc:} the two aliasing charge pairs have equal amplitude products and opposite energy differences, so $\delta_{\rm alias}\propto\cos\tau$ (verified to $1.7\times10^{-18}$). This explains why every failure occurred at $\tau=0$. Finite groups are known to average less completely than $U(1)$ \cite{tran2021}. The skewed label allocation used throughout is an \emph{incomplete} twirl and remains an explicit measurement resource. That recorded labels retain information which the marginal discards is an instance of the Fisher-information chain rule; it has recently been stated for randomized-compiling labels \cite{vyas2026} (preprint).

\section*{Structural null under noise}
\begin{lemma}
Insert channels whose Kraus operators are diagonal in the computational basis (system and clock) anywhere between preparation and the return pulse. The conjugate-class contrast then remains exactly zero.
\end{lemma}
\emph{Proof sketch.} Each Kraus branch preserves the supports $S_n$ and $S_n\cup X_aS_n$, so odd terms still vanish and the transfer path set is unchanged; for the conjugate class it is empty. The direct term carries no common phase on $S_n$. \hfill$\square$

Evidence: phase damping ($p=0.03$) gives $2.56\times10^{-17}$ after the phase layer (stage~19) and $2.1\times10^{-17}$ before the source pulse. Coherent random $R_z$ errors on all eight qubits give $3.4\times10^{-17}$, while the allowed-class contrasts change by up to a factor of three. The lemma does not cover bit-changing errors: a remote $R_x(0.1)$ breaks the null ($2.7\times10^{-5}$), and so do the following exact density-matrix values at $p=0.03$:
\begin{center}\small
\begin{tabular}{@{}l r r r@{}}
\toprule
Channel & Conjugate class & Weakest positive class & $\Lambda_{\rm conjugate}$ (stage-17 budget)\\
\midrule
Ideal & $2.35\times10^{-17}$ & $5.474\times10^{-4}$ & $\simeq0$\\
Preparation amplitude damping & $7.213\times10^{-6}$ & $5.287\times10^{-4}$ & $0.00208$\\
System phase damping & $2.56\times10^{-17}$ & $5.164\times10^{-4}$ & $\simeq0$\\
Symmetric readout flips & $1.651\times10^{-5}$ & $4.397\times10^{-4}$ & $0.00677$\\
\bottomrule
\end{tabular}
\end{center}
The preparation-damping leakage comes from the four system qubits outside the accessed pair. The readout leakage is first order in the flip rate and dominated by the single-bit $q_1$ outcome, so tensored readout mitigation \cite{bravyi2021} should remove it at first order (not tested).

\textbf{Detectability.} $\Lambda$ grows linearly with the number of preparations. Reaching $\Lambda\approx20$ would therefore require about $6.5\times10^{10}$ (readout) or $2.1\times10^{11}$ (damping) preparations at $p=0.03$, and $\Lambda\propto p^2$ at smaller rates. Repetition does not help either. For $k\ge2$ repeated blocks, a source flip in one block followed by another source flip in a later block (and likewise for return flips) returns each $P=0$ state to itself through an out-of-sector intermediate. The conjugate class acquires $2\T_3(0)=14$ such paths, and its ideal contrast is already $5.3\times10^{-4}$ at $k=2$ (stage~27a). Robust phase estimation uses repetition to amplify a rotation angle coherently \cite{kimmel2015}; here, repetition instead destroys the null that defines the background.

\section*{Encoding dependence (stages 24--26)}
With cutoff two and a one-hot local code, the predicted counts become $(5,0,0,2)$, $(19,0,2,2)$ and $(85,0,8,10)$ for $n=2,3,4$. Changing only the phases of invalid multi-hot words alters the same-bit class (by up to $3.42\times10^{-4}$ pointwise) and leaves the other classes exactly invariant. In general, the return probability is invariant when no allowed path passes through an invalid intermediate word, i.e.\ $\Gamma^{\rm inv}_{ab}=\varnothing$. A code-preserving access $R_j(\varepsilon)=\exp(-i\varepsilon H_j/2)$ with $H_j=|0\rangle\langle j|+|j\rangle\langle0|$ removes that sensitivity. It gives the amplitude $\langle r|D_bUD_a|w\rangle+s^2\langle r|H_bUH_a|w\rangle$, with $D_j=I+(c-1)H_j^2$, and the counts $(2,0,0,2)$, $(9,0,2,2)$ and $(37,0,8,10)$. The other-axis/same-sign opening at $n=3$ persists, and the twirl law of Eq.~\eqref{eq:angle} carries over (stage~27a, $5.1\times10^{-19}$). Its contrasts are, however, only $10^{-5}$-scale. Unphysical codewords in unary or one-hot encodings are a standard issue \cite{sawaya2020,lotstedt2025,lamm2020}.

\section*{Independent finite-shot observation}
At $n=3$ and $\varepsilon=\pi/2$, an independently sampled Aer four-test family (22,020,096 preparations; family threshold $0.0125$) detected the same-qubit and other-axis/opposite-sign contrasts ($p=0.00010$ each) but not the conjugate null ($p=0.84162$). The other-axis/same-sign class ($\Lambda=9.90$) missed the threshold ($p=0.02640$). A separately prespecified doubled-budget replication detected it ($\Lambda=19.80$, $p=0.00010$). The two families are not pooled.

\section*{Position relative to prior work}
\textbf{Gauge violation and symmetry protection.} Gauge-breaking terms produce violations growing as $(\lambda t)^2$. Under an energy penalty, they generate an effective in-sector Hamiltonian at second order through virtual out-of-sector excursions \cite{halimeh2020}. Coherent drift out of the physical subspace can be suppressed by random gauge transformations, which act only on errors that dwell outside the sector while a transformation is applied \cite{lamm2020}. Interleaved symmetry transformations average coherent errors \cite{tran2021}, and ancilla-controlled random gauge transformations serve as a Zeno projection \cite{ball2024}. Gauss-law oracles \cite{stryker2019} and symmetry verification \cite{bonetmonroig2018,sagastizabal2019} detect or discard out-of-sector amplitude. In this language, our $A_{ab}$ is the exact two-pulse digital instance of the second-order excursion, the phase layer is the symmetry transformation applied during it, and the label average is the corresponding $\mathbb Z_L$ twirl.

\textbf{Coherent errors and their tailoring.} Coherent-error orientation affects algorithm success \cite{murphy2019}. Randomized compiling tailors coherent errors into Pauli noise \cite{wallman2016,hashim2021}, and syndrome projection in stabilizer codes drives coherent errors towards Pauli-like logical noise \cite{beale2018,bravyi2018}. The sign-twirl result above is a concrete instance of this tailoring for our observable.

\textbf{Reference frames, estimation, and circuits.} Reference frames and superselection \cite{bartlett2007}, quantum Fisher information with and without a phase reference \cite{jarzyna2012}, and imperfect references \cite{safranek2015} form the background of our rank statements. The realizable unitaries of Abelian-symmetric circuits are characterized in \cite{marvian2024}; our pulses deliberately access other sectors, and we predict one return observable. Verified phase estimation uses a return-state test for mitigation \cite{obrien2021}, whereas here the return probability is itself the measured contrast.

\section*{Bounded contribution}
For the declared preparation, two-pulse intervention, clock projection, label distribution and return measurement, we claim:
\begin{enumerate}\itemsep0pt
\item Eqs.~\eqref{eq:amp}--\eqref{eq:angle}: exact, non-perturbative in $\varepsilon$, including the pulse sign and the twirled limit, verified against independent Aer evolution.
\item The support-class path table with its structural conjugate null, the charge-balance reading of the $n=3$ opening, and the invalid-word protection condition $\Gamma^{\rm inv}_{ab}=\varnothing$.
\item The diagonal-noise lemma and computed leakages for selected bit-changing channels.
\item The exact aliasing account of the label resource.
\item Preregistered finite-shot discrimination of the allowed classes from the null (stages 17--18), and of the coherent, sign-flipped and twirled conditions (stage 27b).
\end{enumerate}
We do not claim the general mechanism, a model-independent theorem, universal dimensional selection, practical fault localization, or priority over every possible publication.

\section*{Finite-shot twirl on/off test (stage 27b)}
Before acquisition we froze a protocol (git \texttt{316f0f1}) and the analysis code (\texttt{5f8b8cd}). The family has three conditions---$\sigma=+1$, $\sigma=-1$, and a per-shot random $\sigma$ whose record is discarded---for all four classes, at 1,048,576 shots per label (528,482,304 Aer preparations). A pre-acquisition calibration found nominal 5\% false-positive rates of 0.043--0.051 for the new Wald comparison test and 0.035--0.040 for the weak-null bootstrap.

All six positive predictions passed the family threshold $0.05/6$. Sign dependence gave $\chi^2_3=624.8$, $182.2$ and $547.0$, and the coherent-versus-twirl change gave $77.1$--$127.4$ ($p\le10^{-16}$). The six null predictions, judged at the nominal 0.05, also passed: twirled allowed classes $p=0.66$, $0.28$, $0.22$; conjugate class $p=0.81$, $0.29$, $0.056$. The observed difference vectors had a positive inner product with the exact interference vector for every class, and the 36 standardized residuals against the exact model gave $\sum z^2=43.3$. Exploratorily, re-using the discarded twirl record recovered the interference ($\chi^2_3=92$--$338$). Within this toy circuit, the interference account therefore survives its first falsifiable finite-shot test.

\section*{Reproducibility}
The stage-27a protocol was committed before its code (git \texttt{406b27d}); the results are in \texttt{f3337ee}. For stage 27b, the protocol (\texttt{316f0f1}) and the analysis code (\texttt{5f8b8cd}) both precede acquisition.
\begin{itemize}\itemsep0pt
\item Code (stages 19--27b): \texttt{src/relational\_atlas\_i01/holdout/}; matching JSON in \texttt{results/}.
\item Stage-27b result: \texttt{notes/research\_result27b\_twirl\_shots\_2026-11-13.md}.
\item Literature review: \texttt{notes/literature\_review\_2026-11-13.md}.
\item Stage-27a result: \texttt{notes/research\_result27a\_exact\_design\_2026-11-13.md}.
\item Regression: \texttt{python scripts/test.py} (171 tests).
\end{itemize}

\begin{thebibliography}{99}\small\itemsep0pt
\bibitem{halimeh2020} J. C. Halimeh and P. Hauke, ``Reliability of lattice gauge theories,'' \emph{Phys. Rev. Lett.} \textbf{125}, 030503 (2020).
\bibitem{lamm2020} H. Lamm, S. Lawrence, and Y. Yamauchi, ``Suppressing coherent gauge drift in quantum simulations,'' arXiv:2005.12688 (2020).
\bibitem{tran2021} M. C. Tran, Y. Su, D. Carney, and J. M. Taylor, ``Faster digital quantum simulation by symmetry protection,'' \emph{PRX Quantum} \textbf{2}, 010323 (2021).
\bibitem{ball2024} C. Ball, ``Suppressing gauge drift in quantum simulations with gauge transformations,'' arXiv:2409.04395 (2024).
\bibitem{stryker2019} J. R. Stryker, ``Oracles for Gauss's law on digital quantum computers,'' \emph{Phys. Rev. A} \textbf{99}, 042301 (2019).
\bibitem{bonetmonroig2018} X. Bonet-Monroig, R. Sagastizabal, M. Singh, and T. E. O'Brien, ``Low-cost error mitigation by symmetry verification,'' \emph{Phys. Rev. A} \textbf{98}, 062339 (2018).
\bibitem{sagastizabal2019} R. Sagastizabal \emph{et al.}, ``Experimental error mitigation via symmetry verification in a variational quantum eigensolver,'' \emph{Phys. Rev. A} \textbf{100}, 010302 (2019).
\bibitem{murphy2019} D. C. Murphy and K. R. Brown, ``Controlling error orientation to improve quantum algorithm success rates,'' \emph{Phys. Rev. A} \textbf{99}, 032318 (2019).
\bibitem{wallman2016} J. J. Wallman and J. Emerson, ``Noise tailoring for scalable quantum computation via randomized compiling,'' \emph{Phys. Rev. A} \textbf{94}, 052325 (2016).
\bibitem{hashim2021} A. Hashim \emph{et al.}, ``Randomized compiling for scalable quantum computing on a noisy superconducting quantum processor,'' \emph{Phys. Rev. X} \textbf{11}, 041039 (2021).
\bibitem{beale2018} S. J. Beale, J. J. Wallman, M. Guti\'errez, K. R. Brown, and R. Laflamme, ``Quantum error correction decoheres noise,'' \emph{Phys. Rev. Lett.} \textbf{121}, 190501 (2018).
\bibitem{bravyi2018} S. Bravyi, M. Englbrecht, R. K\"onig, and N. Peard, ``Correcting coherent errors with surface codes,'' \emph{npj Quantum Inf.} \textbf{4}, 55 (2018).
\bibitem{vyas2026} K. Vyas, ``Label and recover coherent errors: randomized compiling does not destroy coherent-error information,'' arXiv:2607.26756 (2026), preprint.
\bibitem{kimmel2015} S. Kimmel, G. H. Low, and T. J. Yoder, ``Robust calibration of a universal single-qubit gate set via robust phase estimation,'' \emph{Phys. Rev. A} \textbf{92}, 062315 (2015).
\bibitem{bravyi2021} S. Bravyi, S. Sheldon, A. Kandala, D. C. Mckay, and J. M. Gambetta, ``Mitigating measurement errors in multiqubit experiments,'' \emph{Phys. Rev. A} \textbf{103}, 042605 (2021).
\bibitem{bartlett2007} S. D. Bartlett, T. Rudolph, and R. W. Spekkens, ``Reference frames, superselection rules, and quantum information,'' \emph{Rev. Mod. Phys.} \textbf{79}, 555 (2007).
\bibitem{jarzyna2012} M. Jarzyna and R. Demkowicz-Dobrza\'nski, ``Quantum interferometry with and without an external phase reference,'' \emph{Phys. Rev. A} \textbf{85}, 011801(R) (2012).
\bibitem{safranek2015} D. \v{S}afr\'anek, M. Ahmadi, and I. Fuentes, ``Quantum parameter estimation with imperfect reference frames,'' \emph{New J. Phys.} \textbf{17}, 033012 (2015).
\bibitem{marvian2024} I. Marvian, ``Theory of quantum circuits with Abelian symmetries,'' \emph{Phys. Rev. Research} \textbf{6}, 043292 (2024).
\bibitem{obrien2021} T. E. O'Brien \emph{et al.}, ``Error mitigation via verified phase estimation,'' \emph{PRX Quantum} \textbf{2}, 020317 (2021).
\bibitem{sawaya2020} N. P. D. Sawaya \emph{et al.}, ``Resource-efficient digital quantum simulation of $d$-level systems for photonic, vibrational, and spin-$s$ Hamiltonians,'' \emph{npj Quantum Inf.} \textbf{6}, 49 (2020).
\bibitem{lotstedt2025} E. L\"otstedt and K. Yamanouchi, ``Comparison of encoding schemes for quantum computing of $S>1/2$ spin chains,'' arXiv:2502.18838 (2025).
\bibitem{oeis} OEIS Foundation, sequences A002426 (central trinomial coefficients) and A005717, \url{https://oeis.org}.
\end{thebibliography}
\end{document}
